\begin{textsample}{NEG Dim 4 – 2016\_c5 – Score 1.00 – 2016\_c5/enade\_2016\_med\_d7a7806f6b.txt}  \label{ex:f4_neg_250}
CAPÍTULO 5 – DISTRIBUIÇÃO DOS CONCEITOS 5.1 - PANORAMA NACIONAL DA DISTRIBUIÇÃO DOS CONCEITOS O Conceito Enade é um indicador de qualidade que avalia o desempenho dos estudantes a partir dos resultados do Enade.
Os valores possíveis vão de 1 ( pior situação ) a 5 ( melhor situação ).
As Unidades de Observação com apenas um ou sem nenhum concluinte participante não obtêm o Conceito Enade, ficando Sem Conceito ( SC ).
A Tabela 5.1 apresenta a quantidade e distribuição de cursos de Medicina participantes do Enade/2016, por faixa de conceito e Grande Região.
Enfatiza -se, mais uma vez, que a diferença entre os cursos tabulados neste capítulo e no capítulo 2 corresponde aos cursos sem conceito, ou seja, àqueles, em princípio, sem alunos concluintes que houvessem participado da prova ou àqueles que tivessem tão somente um aluno concluinte19.
Observando -se os dados da Tabela 5.1, nota -se que, dos 177 cursos participantes, 77 ( 43,5 \% ) classificaram -se com conceito 3, o valor modal.
Esse foi também o conceito modal em quase todas regiões : Norte ( 31,3 \% ), Nordeste ( 48,7 \% ), Sudeste ( 46,8 \% ) e Centro-Oeste ( 41,7 \% ).
Na região Sul, a moda foi o conceito 4, atribuído a 58,1 \% dos cursos da região.
O conceito 4 foi o segundo mais frequente em nível nacional ( 33,3 \%, correspondendo a 59 cursos ) e o conceito 2, o terceiro ( 16,9 \%, correspondendo a 30 cursos ).
Houve, ainda, três cursos ( 1,7 \% ) que receberam conceito 5 e oito cursos ( 4,5 \% ) que receberam conceito 1.
Dos 177 cursos de Medicina, nenhum curso ficou sem conceito ( SC ). 19 Estes cursos com somente um concluinte não foram considerados no capítulo 2.
A região Norte participou com 16 cursos ou 9,0 \% do total nacional.
Desses, cinco cursos receberam o conceito 3, o que equivale a 31,3 \% do total regional, além de ser o conceito modal na região, como já comentado.
A três cursos ( 18,8 \% ) atribuiu -se conceito 1, e a quatro cursos cada ( 25,0 \% cada ), os conceitos 2 e 4.
Nenhum curso ficou com conceito 5 nessa região.
A região Nordeste participou com 39 cursos, ou 22,0 \% do total nacional.
Nessa região, 19 cursos ( 48,7 \% em termos regionais ) obtiveram conceito 3, o conceito modal para a região, como já comentado.
Os conceitos 2 e 4 foram atribuídos, respectivamente, a sete e a 12 cursos ( 17,9 \% e 30,8 \%, respectivamente ).
Ainda nessa região, um curso ( 2,6 \% ) foi avaliado com o conceito 5.
Nenhum curso ficou com conceito 1.
Dos 79 cursos participantes da região Sudeste ( 44,6 \% do total nacional ), o conceito 3 foi o conceito modal, como já anteriormente assinalado, tendo sido obtido por 37 cursos ( 46,8 \% ).
O conceito 1 foi atribuído a cinco cursos ( 6,3 \% ) e o conceito 2, a 15 cursos ( 19,0 \% ).
Receberam os conceitos 4 e 5, respectivamente, 20 cursos ( 25,3 \% ) e dois cursos ( 2,5 \% ).
Os 31 cursos da região Sul corresponderam a 17,5 \% do total nacional.
Como já indicado, houve predominância do conceito 4, atribuído a 58,1 \% dos cursos da região ( conceito modal ), ou seja, atribuído a 18 dos 31 cursos participantes na região Sul.
O conceito 2 foi atribuído a dois cursos ( 6,5 \% ) e o conceito 3, a 11 cursos ( 35,5 \% ).
Nenhum curso recebeu conceito 1 ou 5 na região Sul.
Já dos 12 cursos participantes na região Centro-Oeste ( 6,8 \% do total nacional ), nota -se que tais cursos apresentaram distribuição bimodal, uma vez que as maiores frequências se concentraram em dois conceitos, os conceitos 3 e 4, atribuídos, cada qual, a cinco cursos ( 41,7 \% cada ).
Os dois cursos restantes foram atribuídos ao conceito 2 e nenhum curso ficou com conceito 1 ou 5.
O Gráfico 5.1 apresenta a distribuição do Conceito Enade segundo Grande Região : poligonais mais à esquerda representam Grandes Regiões com pior distribuição, e poligonais mais à direita, Grandes Regiões com melhores conceitos.
A região Norte ( linha azul ) apresenta os cursos com conceitos mais baixos, ao passo que a região Sul ( linha roxa ) apresenta a distribuição com valores maiores.
Cursos em IES situadas nas regiões Nordeste ( linha verde ), Sudeste ( linha vermelha ) e Centro-Oeste ( linha preta ) apresentam uma situação intermediária. 5.2 - CONCEITOS POR CATEGORIA ADMINISTRATIVA E GRANDE REGIÃO A Tabela 5.2 apresenta a distribuição dos cursos participantes do Enade/2016 de Medicina, por Categoria Administrativa, de acordo com os conceitos por eles alcançados, segundo as Grandes Regiões brasileiras.
Dos 177 cursos participantes, 75 ( 42,4 \% ) eram ministrados em Instituições Públicas e 102 ( 57,6 \% ), em Instituições Privadas.
De acordo com as informações da Tabela 5.2, em termos nacionais, entre Instituições Públicas e Privadas, observa -se que os três cursos avaliados com conceito 5 eram oferecidos em IES Públicas.
Além disso, em termos nacionais, nenhum curso ficou sem conceito ( SC ).
Dos 75 cursos participantes de IES Públicas, o conceito 4 foi o valor modal, atribuído a 40 cursos ( 53,3 \% ).
Entre os demais cursos de Instituições Públicas participantes, seis cursos ( 8,0 \% da categoria ) obtiveram conceito 2, 26 cursos ( 34,7 \% ) receberam conceito 3 e, como já comentado, três cursos ( 4,0 \% ) receberam conceito 5.
Nesta categoria, nenhum curso obteve conceito 1 ou ficaram sem conceito ( SC ).
Na rede privada, o conceito modal foi 3, atribuído a 51 cursos, o correspondente a 50,0 \% dos 102 cursos da categoria.
Entre os demais cursos participantes, oito ( 7,8 \% ) receberam conceito 1, e 24 ( 23,5 \% ), conceito 2.
O conceito 4 foi atribuído a 19 cursos ( 18,6 \% ).
Nesta Categoria Administrativa, nenhum curso ficou com conceito 5 ou sem conceito ( SC ).
Na análise por região, observa -se que, na região Norte, as Instituições Privadas participaram com sete cursos ( 43,8 \% do total regional ), dos quais nenhum recebeu conceito 5.
O conceito modal para as Instituições Privadas na região foi 1, atribuído a três cursos, correspondendo a 42,9 \% dessa categoria na região.
Dois cursos foram avaliados com conceito 2, e outros dois com conceito 3 e 4 ( um curso, cada ).
As Instituições Públicas participaram com nove cursos na região Norte ( 56,3 \% do total regional ), dos quais quatro obtiveram conceito 3, o valor modal.
O conceito 2 foi atribuído a dois cursos, e o conceito 4 a três cursos.
Nenhum curso ficou com os conceitos 1 ou 5 nessa combinação de categoria e região.
Na região Nordeste, a rede privada concentrou 14 dos 39 cursos participantes, o equivalente a 35,9 \% do total da região, a menor das proporções desta categoria entre as regiões.
Dentre os cursos oferecidos por IES Privadas no Nordeste, oito cursos foram avaliados com conceito 3, o conceito modal.
Quatro cursos obtiveram conceito 2 e dois cursos receberam conceito 4.
Nessa combinação de categoria e região, nenhum curso com os conceitos 1 ou 5.
As Instituições Públicas da região Nordeste participaram com 25 cursos ( 64,1 \% do total da região ), a maior das proporções desta categoria entre as regiões.
Desses, 11 obtiveram conceito 3, o valor modal, seguidos por outros dez cursos, que receberam conceito 4.
Os conceitos 2 e 5 foram atribuídos, respectivamente, a três e a um curso.
Nenhum curso de IES Pública recebeu conceito 1 nessa região.
Na região Sudeste, a proporção de cursos da rede privada, 72,2 \%, foi mais elevada do que a registrada nas demais regiões brasileiras, correspondendo a 57 dos 79 cursos participantes.
Nessa categoria e região, o conceito modal foi 3, atribuído a 30 cursos.
Os demais cursos foram avaliados com conceito 1 ( cinco cursos ), conceito 2 ( 15 cursos ), conceito 4 ( sete cursos ).
Entre os 22 cursos oferecidos por Instituições Públicas na região Sudeste, o conceito modal foi 4, atribuído a 13 cursos.
Os demais cursos receberam os conceitos 3 ( sete cursos ) e 5 ( dois cursos ).
Na região Sudeste, nesta categoria, nenhum curso obteve os demais conceitos.
As Instituições Privadas concentraram 18 dos 31 cursos participantes da região Sul, 58,1 \% do total regional.
Desses, nove obtiveram conceito 3, o conceito modal.
Nesta combinação de Categoria Administrativa e Grande Região, os conceitos 2 e 4 foram atribuídos a um e a oito cursos, respectivamente.
Nenhum curso recebeu os conceitos 1 ou 5 nessa combinação de categoria e região.
As Instituições Públicas na região Sul participaram com 13 cursos, 41,9 \% dos cursos da região.
O conceito modal foi 4, atribuído a dez cursos.
Aos demais cursos foram atribuídos os conceitos 3 ( dois cursos ) e 2 ( um curso ).
Nenhum curso da região Sul oferecido por IES Pública ficou com os conceitos 1 ou 5.
Na região Centro-Oeste, metade dos cursos participantes ( seis dos 12 ) eram de Instituições Privadas ( 50,0 \% em termos regionais ).
Desses, três receberam conceito 3, o conceito modal.
Entre os demais, dois receberam conceito 2 e um ficou com conceito Na região Centro-Oeste, metade dos cursos participantes ( seis dos 12 ) eram de Instituições Privadas ( 50,0 \% em termos regionais ).
Desses, três receberam conceito 3, o conceito modal.
Entre os demais, dois receberam conceito 2 e um ficou com conceito 4.
Nenhum curso ficou com os conceito 1 ou 5.
Quanto a outra metade dos cursos oferecidos por Instituições Públicas na região Centro-Oeste ( 50,0 \% do total regional ), nota -se que quatro cursos apresentaram conceito 4 ( conceito modal ) e dois ficaram com conceito 3.
Nenhum curso ficou com os demais conceitos nessa combinação de categoria e região.
O Gráfico 5.2 apresenta a distribuição do Conceito Enade segundo a Categoria Administrativa da IES.
Os cursos em IES Públicas ( linha azul ) apresentam uma poligonal mais à direita do que os em IES Privadas ( linha verde ) e, consequentemente, uma distribuição de Conceitos Enade com valores maiores. 5.3 - CONCEITOS POR ORGANIZAÇÃO ACADÊMICA E GRANDE REGIÃO Na Tabela 5.3, encontra -se a distribuição dos conceitos atribuídos aos cursos participantes do Enade/2016 na Área de Medicina, por Organização Acadêmica, segundo as Grandes Regiões brasileiras.
Dos 177 cursos de Medicina participantes, 121 eram oferecidos em Universidades, 18 em Centros Universitários e 38 em Faculdades.
Esta distribuição corresponde a, respectivamente, 68,4 \%, 10,2 \% e 21,5 \% dos cursos.
Além disso, no Brasil como um todo, nenhum curso ficou sem conceito ( SC ).
De acordo com os dados apresentados, dos três cursos avaliados com conceito 5, dois eram vinculados a Universidades.
Os cursos vinculados a esse tipo de Organização Acadêmica tiveram o conceito 3 como conceito modal, atribuído a 50 cursos ( 41,3 \% ).
Os demais cursos vinculados a Universidades receberam os conceitos 1 ( cinco cursos ), 2 ( 16 cursos ), 4 ( 48 cursos ) e conceito 5 ( dois cursos, como já mencionado ).
Entre os cursos em Centros Universitários, o conceito modal foi também 3, atribuído a oito cursos ( 44,4 \% ).
Os outros cursos receberam os conceitos 1 ( dois cursos ), 2 ( sete cursos ) e 4 ( um curso ).
Nenhum curso obteve conceito 5.
Dos 38 cursos mantidos por Faculdades, 19 ( 50,0 \% ) receberam conceito 3, o conceito modal.
O segundo conceito mais frequentemente atribuído aos cursos oferecidos por essa categoria de IES foi o conceito 4, recebido por dez cursos.
Na sequência, sete cursos obtiveram conceito 2 e dois cursos ficaram com os conceitos 1 e 5 ( um curso, cada ).
Considerando -se separadamente as regiões brasileiras, verifica -se que, na região Norte, as Universidades concentraram nove dos 16 cursos participantes.
Quatro cursos receberam conceito 3, conceito modal.
Os demais cursos receberam os conceitos 1 e 2 ( um curso, cada ) e 4 ( três cursos ).
Nenhum curso ficou recebeu conceito 5 nessa combinação de organização e região.
Considerando -se separadamente as regiões brasileiras, verifica -se que, na região Norte, as Universidades concentraram nove dos 16 cursos participantes.
Quatro cursos receberam conceito 3, conceito modal.
Os demais cursos receberam os conceitos 1 e 2 ( um curso, cada ) e 4 ( três cursos ).
Nenhum curso ficou recebeu conceito 5 nessa combinação de organização e região.
Os Centros Universitários da região Norte foram representados por três cursos, aos quais foram atribuídos os conceitos 1, 2 e 4 ( um curso, cada ).
Nenhum curso obteve o conceito 3 ou 5.
As Faculdades participaram com quatro cursos na região Norte que receberam os conceitos 1 e 3 ( um curso, dada ) e 2 ( dois cursos, conceito modal ).
Nenhuma Faculdade da região ficou com conceito 4 ou 5.
Na região Nordeste, as Universidades participaram com 29 dos 39 cursos da Área de Medicina oferecidos na região.
O conceito modal foi 3, atribuído a 13 cursos, enquanto o conceito 5 foi atribuído a um curso.
Obtiveram os conceitos 2 e 4, respectivamente, quatro e 11 cursos.
Nenhum curso oferecido por Universidades no Nordeste recebeu conceito 1.
Os Centros Universitários contaram com dois cursos participante na região Nordeste.
Esses obtiveram conceito 3.
As Faculdades foram representadas por oito cursos na região Nordeste, quatro dos quais obtiveram o conceito modal 3.
O conceito 2 foi atribuído a três cursos e o conceito 4, a um curso.
Na região Nordeste, nenhum curso oferecido por Faculdades recebeu conceito 1 ou 5.
Na região Sudeste, as Universidades concentraram 46 dos 79 cursos de Medicina da região.
Entre os cursos em Universidades, nessa região, o conceito modal foi 3, atribuído a 19 cursos.
Os demais cursos receberam os conceitos 1 ( quatro cursos ), 2 ( nove cursos ), 4 ( 13 cursos ) e 5 ( um curso ).
Os Centros Universitários participaram com dez cursos na região Sudeste, dos quais cinco obtiveram o conceito modal 2.
O conceito 1 foi atribuído a um único curso, e o conceito 3, a quatro cursos.
Nenhum curso ficou com conceito 4 ou 5 nessa combinação de Organização Acadêmica e região.
Já as Faculdades foram representadas por 23 cursos na região Sudeste, 14 dos quais ficaram com conceito 3, o conceito modal.
Aos demais cursos foram atribuídos os conceitos 2 e 5 ( um curso, cada ) e 4 ( sete cursos ).
Dos 31 cursos da região Sul, 28 eram de Universidades, para os quais o conceito modal foi 4, atribuído a 17 cursos.
Nesse tipo de organização, nenhum curso ficou com conceito 1 ou 5.
Os demais cursos receberam os conceitos 2 ( um curso ) e 4 ( 17 cursos ).
Os Centros Universitários tiveram dois cursos participantes na região Sul, esses dois cursos obtiveram os conceitos 2 e 3 ( um curso, cada ).
Já as Faculdades participaram com um único curso na região Sul, esse curso recebeu o conceito 4.
Na região Centro-Oeste, nove dos 12 cursos eram de Universidades.
Nesse tipo de organização, nenhum curso recebeu conceito 1 ou 5.
O conceito 2 foi atribuído a um curso e o conceito 3 e 4, conceitos modais, a quatro cursos, cada.
Os Centros Universitários da região Centro-Oeste contaram com somente um curso participante, que obteve o conceito 3.
As Faculdades ofereceram dois cursos participantes, dos quais um foi \textbf{alocado} ao conceito 2 e outro ao conceito 4.
O Gráfico 5.3 apresenta a distribuição do Conceito Enade segundo a Organização Acadêmica da IES.
Os cursos em Universidades ( linha azul ) apresentam uma poligonal mais à direita do que as demais, denotando uma melhor distribuição de Conceitos Enade.
No outro extremo, temos os cursos oferecidos em Centros Universitários ( linha verde ) com a poligonal mais à esquerda e, consequentemente, uma distribuição de Conceitos Enade com menores valores.
IES de Faculdades ( linha vermelha ) apresentam uma situação intermediária.

% matched lemmas: alocar
\end{textsample}
